\clearpage\SourceLeaf{316}

\begin{center}
[II.]

[DIE PHYSISCHE MÖGLICHKEIT DER DATEN IN POTHENOTS AUFGABE.]

[1.]
\end{center}

\begin{center}
\textsc{Gauss} an \textsc{Schumacher}. Göttingen, 13. April 1836.
\end{center}

\indent \mbox{. . . . . . .} Bei dieser Gelegenheit will ich auch eine andere Aufgabe erwähnen, nämlich die sogenannte \textsc{Pothenot}sche Aufgabe. Sind drei Punkte $A$, $B$, $C$ gegeben und zugleich die Winkel, die $XA$, $XB$, $XC$ miteinander bilden (natürlich einschließlich des \emph{sens}, d.\,h. der Angabe, ob die Winkel von der linken zur rechten [Seite] oder umgekehrt gezählt werden), so ist $X$ (im Allgemeinen gesprochen, d.\,h. mit Ausnahme des Falls, in dem $A$, $B$, $C$, $X$ auf einem Kreis liegen) \textls[120]{bestimmt}. Beruhen die Daten auf einem wirklichen konkreten Fall, so versteht es sich von selbst, dass die Aufgabe, weil sie wirklich vorliegt, auch möglich ist. Wenn Ihnen aber jemand (um Sie irrezuführen) bloß die Zahlen vorgibt, kann sich im Nachhinein zeigen, dass die Aufgabe physisch unmöglich ist und dass Sie anstelle einer der drei Richtungen die entgegengesetzte erhalten (die sich um $180^0$ von der vorgegebenen unterscheidet).

Die Frage ist nun, wie man \textls[120]{am einfachsten} die Bedingung der physischen Möglichkeit durch eine Gleichung \textls[120]{zwischen den Daten} ausdrücken kann. Da ich mich darüber wunderte, dass bei einer Aufgabe, über die so viel geschrieben wurde, doch alle, soweit ich weiß, einen so wesentlichen Umstand völlig übersehen haben, habe ich die Lösung oft in meinen Vorlesungen vorgetragen. Aber ganz so leicht ist sie wohl doch nicht zu finden. Wenigstens habe ich sie einmal einem sonst sehr fähigen jungen Mann, dem ich die Lösung \textls[120]{nicht} vorgetragen hatte, als Thema einer Doktorarbeit vorgeschlagen. Obwohl er sich lange damit abmühte, konnte er die Lösung nicht finden. \mbox{. . . . . . .}
